% ===================================================================
%  図表第 9 号 — 山。— 温泉町。— 森。— 狩。
%
%  Ceci n'est PAS une transcription : c'est une traduction, faite en
%  2026, en japonais d'aujourd'hui. Voir texto/ja/10-zuhyo-01.tex
%  pour la consigne, qui vaut pour les seize fichiers.
% ===================================================================

%%K t09-apar-1 apar
\VUsaut{4.0mm}
\VUtitre{15.0pt}{60}{図表第 9 号}
\VUsaut{3.0mm}

%%K t09-tit-1 sub
\VUcentre{12.0pt}{0}{\textit{第一場。}}
\VUsaut{3.0mm}

%%K t09-tit-2 sub
\VUpk{11.4pt}{40}{山。--- 温泉町。}
\VUsaut{3.0mm}

%%K t09-c1-01-1 p
1. --- 私はプラッツに来て一週間になる。ピレネエのあの見事な\VUgras{山脈}
\textsuperscript{(1)} の前に立ったときの驚きは、とうてい言い表せない。その
\VUgras{尾根} \textsuperscript{(2)} が青空に刻む輪郭は、巨大な鋸のようである。

%%K t09-c1-02-1 p
2. --- ピレネエの\VUgras{峰} \textsuperscript{(3)} がこれほどの\VUgras{高さ}に
達しうるとは思わなかった。あの壮麗な\VUgras{氷河} \textsuperscript{(5)} と雪を
かぶった\VUgras{尖峰} \textsuperscript{(4)} に私は見とれてしまった。恐ろしい
\VUgras{雪崩}はその上から落ちてくるのである。

%%K t09-tit-3 sub
\VUsaut{3.0mm}
\VUpk{10.2pt}{40}{山の景色。}
\VUsaut{3.0mm}

%%K t09-c1-03-1 p
3. --- 着いた翌日、私はプラッツの近くのとうに消えた\VUgras{火山}
\textsuperscript{(6)} へ行った。\VUgras{噴火口}の裸で痩せた縁には、幾百年も前に
\VUgras{山の斜面} \textsuperscript{(7)} を流れ下った\VUgras{熔岩}の跡がいまでも
はっきり見える。ある\VUgras{斜面}で私はその熔岩の欠片を幾つか見つけることが
できた。

%%K t09-c1-04-1 p
4. --- プラッツは国境にある。あの\VUgras{要塞} \textsuperscript{(8)} は、
フランスとスペインとのあいだの\VUgras{峠} \textsuperscript{(27)} を守っていて、
ライン河のほとりのあの古い城によく似ている。それらは自分の岩の上から\VUgras{獲物}をうかがって
いるように見えるのである。

%%K t09-c1-05-1 p
5. --- 巨大な\VUgras{鷲} \textsuperscript{(9)} がしばしば私の目を引く。その巣は
\VUgras{砦} \textsuperscript{(8)} から遠くない。この\VUgras{猛禽}は力強い
\VUgras{嘴}をもち、しばしば雛のために仔羊や仔山羊を自分の\VUgras{爪}に掴んで
運ぶ。その\VUgras{翼}はたいそう大きいので、重い荷を負ってもながく滑ることが
できる。

%%K t09-tit-4 sub
\VUsaut{3.0mm}
\VUpk{10.2pt}{40}{歩く人。}
\VUsaut{3.0mm}

%%K t09-c1-06-1 p
6. --- ここで一番心地よい道行きはコの\VUgras{滝} \textsuperscript{(10)} まで行く
ことである。\VUgras{落ちる水}は渦を巻いて下り、やがて美しい\VUgras{小川}と
なって町の西の\VUgras{樅の林} \textsuperscript{(11)} に消えてゆく。我々はこの林を
抜けて\VUgras{測候所} \textsuperscript{(12)} まで登り、\VUgras{望楼}の上から
この地方の\VUgras{全景}をひと目に収めた。\VUgras{眺め}はきわめてよく、
\VUgras{自然}が手元の宝をことごとくこの土地に注いだかのようである。

%%K t09-c1-07-1 p
7. --- 下りには\VUgras{綱の鉄道} \textsuperscript{(13)} に乗り、小さな\VUgras{駅}
で降りた。その駅は昔プラッツの領主が住んだ古い\VUgras{領主の城}の
\VUgras{廃墟} \textsuperscript{(14)} のそばである。

%%K t09-c1-08-1 p
8. --- 気の毒な城よ。足もとのあの小ぎれいな\VUgras{温泉町} \textsuperscript{(15)}
が、いまや流行の\VUgras{湯治場}となったのを見て、何を思うことであろう。

%%K t09-c1-08-2 p
ここの\VUgras{気候}はたいそう快く、\VUgras{鉱泉}はたいそう効くので、
\VUgras{療養所} \textsuperscript{(17)} で受ける一\VUgras{療程}はまったくの楽しみで
ある。

%%K t09-tit-5 sub
\VUsaut{3.0mm}
\VUpk{10.2pt}{40}{快い温泉町。}
\VUsaut{3.0mm}

%%K t09-c1-09-1 p
9. --- まことに、居心地よく過ごせるようあらゆる手が尽くされている。鉄道を
通すために大きな\VUgras{高架橋} \textsuperscript{(16)} が架けられた。\VUgras{湯の館}
の\VUgras{庭} \textsuperscript{(18)} は洒落たしつらえで、そこでは\VUgras{音楽会}も
開かれる。美しい\VUgras{別荘} \textsuperscript{(19)} が幾つも\VUgras{賭博場}
\textsuperscript{(21)} のまわりに建ち、よく手入れした\VUgras{街道}
\textsuperscript{(22)} が行き来をきわめてたやすくしている。

%%K t09-c1-10-1 p
10. --- ただ一つの\VUgras{土手} \textsuperscript{(23)} が\VUgras{道}
\textsuperscript{(20)} と\VUgras{谷} \textsuperscript{(25)} そのものとを隔てている。
谷の底には美しい小さな\VUgras{湖} \textsuperscript{(26)} が横たわり、その水が澄んだ
\VUgras{小川} \textsuperscript{(24)} を養っている。

%%K t09-c1-11-1 p
11. --- 谷の向う側では鉄の\VUgras{鉱山} \textsuperscript{(28)} を掘っている。私は
何度も\VUgras{坑夫}が\VUgras{立坑}を下るのを見に行き、\VUgras{横坑}にも入った。
彼らはそこで一日じゅう\VUgras{金属}を含む\VUgras{鉱石}を掘り出している。

%%K t09-c1-12-1 p
12. --- 鉱山のそばに大きな\VUgras{製錬所}が建てられた。掘り出した富をその場で
用いるためである。\VUgras{溶鉱炉} \textsuperscript{(29)} はここに多い\VUgras{石炭}
で焚かれ、\VUgras{銑鉄}を早く安く得ることを許す。

%%K t09-c1-13-1 p
13. --- 鉱山から遠くない所に\VUgras{石切場} \textsuperscript{(30)} を開いている。
老いた\VUgras{石切}が\VUgras{鶴嘴}で割り出すのは\VUgras{花崗岩}でも、
\VUgras{斑岩}でも、ましてや\VUgras{砂岩}でもなく、家を建てるためのただの
\VUgras{建築石}である。

%%K t09-c1-14-1 p
14. --- この土地の最も美しい飾りの一つは\VUgras{樅} \textsuperscript{(31)} である。
どこにでもある --- \VUgras{峡} \textsuperscript{(32)} の底にも、\VUgras{急流}
\textsuperscript{(33)} のほとりにも、\VUgras{深い淵} \textsuperscript{(34)} や崖の
そばにも --- そしてその細い針の葉が緑の一筆を添える。

%%K t09-tit-6 sub
\VUsaut{3.0mm}
\VUpk{10.2pt}{40}{歩くこと。}
\VUsaut{3.0mm}

%%K t09-c1-15-1 p
15. --- 私は休みを利して長く歩いている。山を巻いてゆく\VUgras{小道}
\textsuperscript{(35)} は、それと気づかぬうちに幾\VUgras{キロメートル}も、しばしば
幾\VUgras{里}も歩かせてしまう。

%%K t09-c1-15-2 p
\VUgras{里程標} \textsuperscript{(36)} と\VUgras{道しるべ} \textsuperscript{(37)} が
小道に印をつけている。その道でしばしば若い\VUgras{山の男} \textsuperscript{(38)} に
会う。腰に\VUgras{頭陀袋} \textsuperscript{(40)} を掛け、手に\VUgras{土竜鼠}
\textsuperscript{(39)} を抱えている。あるいは\VUgras{流浪の民} \textsuperscript{(41)}
に会う。その\VUgras{毛皮の帽子} \textsuperscript{(42)} は破れた古い\VUgras{外套}
\textsuperscript{(43)} と妙に釣り合わない。彼は仕込んだ\VUgras{熊}
\textsuperscript{(44)} を\VUgras{鎖}で引いてゆく。

%%K t09-tit-7 sub
\VUpk{10.2pt}{40}{山登り。}
\VUsaut{3.0mm}

%%K t09-c1-16-1 p
16. --- 私はほとんど毎日\VUgras{案内人} \textsuperscript{(48)} と\VUgras{山登り}に
行く。背に\VUgras{背嚢} \textsuperscript{(47)} を負い、手に\VUgras{登山杖}を持ち、
まことの\VUgras{旅人} \textsuperscript{(46)} のように\VUgras{岩} \textsuperscript{(45)}
に挑むときは、たいそう誇らしい。

%%K t09-c1-17-1 p
17. --- 山の頂から見ると、平地の人は蟻ほどの大きさもない。\VUgras{羊の群} \textsuperscript{(50)}
は子供の玩具のようである。群は\VUgras{牧場} \textsuperscript{(49)} で草を食んで
いる。一本の\VUgras{松} \textsuperscript{(51)} も、木の\VUgras{山小屋}
\textsuperscript{(52)} さえも、緑の上の小さな点にすぎない。

%%K t09-c1-18-1 p
18. --- こうした道行きでは\VUgras{山道} \textsuperscript{(53)} でしばしば
\VUgras{羚羊撃} \textsuperscript{(54)} に出会う。いま仕留めた獣を肩に担いで
いる。

%%K t09-c1-19-1 p
19. --- 私は標本帖にたいそう珍しい\VUgras{羊歯} \textsuperscript{(55)} の葉と、
美しい桃色の\VUgras{石南} \textsuperscript{(57)} の枝を収めた。前の散歩のとき、
小さな\VUgras{洞} \textsuperscript{(56)} の口で摘んだものである。

%%K t09-tit-8 sub
\VUsaut{3.0mm}
\VUcentre{12.0pt}{0}{\textit{第二場。}}
\VUsaut{3.0mm}

%%K t09-tit-9 sub
\VUpk{11.4pt}{40}{森。--- 狩。}
\VUsaut{3.0mm}

%%K t09-c2-01-1 p
1. --- 昨日私は森でこの上なく楽しい一日を過ごした。そこに住む獣や
\VUgras{虫} \textsuperscript{(5)} の忙しい暮しが、私にはたいそう面白かった。

%%K t09-c2-01-2 p
\VUgras{蜘蛛} \textsuperscript{(1)} は二本の枝のあいだに\VUgras{網}
\textsuperscript{(2)} を張り、通りかかる\VUgras{蠅} \textsuperscript{(3)} を捕える。
\VUgras{蝸牛} \textsuperscript{(4)} は殻を負って重たげに這う。鍬形は獲物を探し、
働き者の蟻は\VUgras{蟻塚} \textsuperscript{(6)} のそばで精を出す --- この小さな者
たちが行き来し、驚くほどの熱心さで働いていた。

%%K t09-c2-02-1 p
2. --- 一匹の\VUgras{蛇} \textsuperscript{(7)} は、初め\VUgras{蝮}かと思ったが、
じつは咬まぬただの\VUgras{水蛇}であった。木の幹の下にとぐろを巻き、ときおり
細い小さな頭を出す。その頭はいつでも咬みつきそうに見えた。私の前では大きな
\VUgras{なめくじ} \textsuperscript{(8)} が重たげに這っていた。ここに一面に生える
おいしい小さな\VUgras{野苺} \textsuperscript{(11)} を一つ平らげたばかりである。

%%K t09-c2-03-1 p
3. --- 地は\VUgras{森の花}に敷きつめられていた。\VUgras{桜草}
\textsuperscript{(9)}、\VUgras{蔓日日草} \textsuperscript{(10)}、それに私の知らない
草がたくさん、\VUgras{草の叢} \textsuperscript{(12)} のあいだに咲いていた。

%%K t09-c2-04-1 p
4. --- このあたりでは\VUgras{松露}が\VUgras{茸} \textsuperscript{(14)} と同じ
くらい多い。ある森番の\VUgras{仔豚} \textsuperscript{(13)} が、幾つか見つからぬ
ものかと欲深げに鼻で掘っていた。

%%K t09-tit-10 sub
\VUsaut{3.0mm}
\VUpk{10.2pt}{40}{密猟。}
\VUsaut{3.0mm}

%%K t09-c2-05-1 p
5. --- 私はたいそう愉快な一幕の\VUgras{終り}を見た。自分の\VUgras{短脚の猟犬}
\textsuperscript{(15)} の鼻のおかげで、\VUgras{森番} \textsuperscript{(16)} はちょうど
\VUgras{密猟者} \textsuperscript{(22)} が仕留めた\VUgras{獲物} \textsuperscript{(17)}
を運び去ろうとするところを取り押えたのである。密猟者は捕まるよりは獲物を
捨てる方を選んで逃げ、\VUgras{野兎} \textsuperscript{(18)}、\VUgras{牝鹿}
\textsuperscript{(19)}、\VUgras{麕} \textsuperscript{(20)}、\VUgras{猪}
\textsuperscript{(21)} が草の上に残された。森番にはこの上ない喜びである。

%%K t09-c2-06-1 p
6. --- この森の木の種類の多さには驚かされる。\VUgras{橅} \textsuperscript{(23)} と
\VUgras{白樺} \textsuperscript{(26)}、\VUgras{松脂} \textsuperscript{(27)} を出す
\VUgras{松}、\VUgras{樫} \textsuperscript{(29)}、そのほか多くの木が枝を交わして
いる。

%%K t09-c2-06-2 p
高い幹に絡む\VUgras{蔦} \textsuperscript{(24)} がこの\VUgras{高木の林}
\textsuperscript{(25)} に鮮やかな緑を添え、それが\VUgras{苔} \textsuperscript{(31)} の
淡くやわらかな緑とよく調和している。\VUgras{青啄木鳥} \textsuperscript{(30)} は
その苔の中に虫を探している。

%%K t09-tit-11 sub
\VUsaut{3.0mm}
\VUpk{10.2pt}{40}{森の景色。}
\VUsaut{3.0mm}

%%K t09-c2-07-1 p
7. --- あの老いた樫に私は心を奪われる。太く曲がった\VUgras{根}
\textsuperscript{(32)} が半ば土の外に現れ、いかにも力強い姿を与えている。
\VUgras{林の主} \textsuperscript{(35)} がどうしてそれを倒させ、\VUgras{鋸}
\textsuperscript{(34)} に渡す気になれるのか、私には分からない。その鋸は
\VUgras{木樵} \textsuperscript{(33)} のものである。気の毒な\VUgras{幹} \textsuperscript{(36)} が\VUgras{皮}
\textsuperscript{(37)} を剥がれ、あるいは\VUgras{斧} \textsuperscript{(38)} で割られて
いる --- それは\VUgras{日雇} \textsuperscript{(39)} の斧である --- のを見るのは
辛い。

%%K t09-c2-08-1 p
8. --- そのそばで老いた\VUgras{炭焼} \textsuperscript{(41)} が自分の\VUgras{鋤}で
炭の\VUgras{山} \textsuperscript{(42)} を積んでおり、一人の老婆が\VUgras{枯枝}の
\VUgras{束} \textsuperscript{(40)} を背負って行くところであった。倒された木の枝で
作った束である。

%%K t09-tit-12 sub
\VUsaut{3.0mm}
\VUpk{10.2pt}{40}{狩。}
\VUsaut{3.0mm}

%%K t09-c2-09-1 p
9. --- 私は\VUgras{森番の家} \textsuperscript{(46)} へ休みに行くつもりであったが、
小さな\VUgras{池} \textsuperscript{(43)} まで来て --- そこには美しい\VUgras{白鳥}
\textsuperscript{(45)} が泳いでいた --- 先ごろの\VUgras{洪水} \textsuperscript{(44)}
が道をまったくの\VUgras{沼}にしてしまったのに気づいた。私は回り道をせねばなら
ず、森のはずれの\VUgras{杉} \textsuperscript{(49)}、\VUgras{ポプラ}
\textsuperscript{(48)}、\VUgras{楡} \textsuperscript{(47)} のそばを通るとき、
\VUgras{茂み} \textsuperscript{(54)} から見事な\VUgras{牡鹿} \textsuperscript{(50)}
が飛び出すのを見た。その後ろには容赦のない\VUgras{猟犬の群} \textsuperscript{(51)}
が迫っており、\VUgras{角笛} \textsuperscript{(53)} がそれを煽っていた。その角笛は
\VUgras{馬上の狩人} \textsuperscript{(52)} のものである。哀れな獣はもう力尽きて
いた。私から数歩の所で息絶えて倒れた。

%%K t09-c2-10-1 p
10. --- 森番の子が\VUgras{狩}の物音に引かれて出てきて、我々二人はまた森へ
入った。その子は老いた樫の枝に\VUgras{鵲} \textsuperscript{(56)} を一羽見つけて
いた。その巣は\VUgras{葉むら} \textsuperscript{(58)} に隠れているに違いないと思い、
取りたかったのである。

%%K t09-c2-10-2 p
\VUgras{栗鼠} \textsuperscript{(61)} のように身軽に、その子は\VUgras{枝}
\textsuperscript{(55)} から枝へと登ったが、何も見つからず、ただ一羽の老いた
\VUgras{烏} \textsuperscript{(60)} を驚かせただけであった。烏は\VUgras{大枝}
\textsuperscript{(57)} に止まっていたのである。
\VUsaut{4.0mm}
\VUfilet{22mm}
